\section{Computable centered ground energy with finite null constraints}
\label{vfourC:section}

This appendix integrates the conditional enclosure construction of
\cite{vfourComputable}.  Its object is the bottom of the \emph{actual complete}
centered form, not the bottom of a fixed prime head or a finite trial matrix.
Throughout this appendix, as in the main text,
\[
 \ip{u}{v}_{\Hh}=\int_{\R}u(x)\overline{v(x)}\,\frac{dx}{W(x)}
\]
is linear in its first argument.  The Fourier transform is the ordinary unitary
transform specified in the common notation.  In particular, matrix entries and
Fourier features below use this convention, rather than the opposite convention
in the standalone source.

\subsection{Realization input and the enclosure theorem}
\label{vfourC:subsec:input}

We use the complete forms \eqref{a:eq:physicalgamma}--\eqref{a:eq:qzero} on
$X$ of \eqref{a:eq:X}.  Write $L:=B=T_\Gamma-D_b$ for the form-defined weighted
representative in Section~\ref{sec:weighted-physical}, and $L_c$ for its
restriction to $\Hh_c$.  The analytic realization inputs are the following:
$X$ is the closed form domain; the complete prime correction is bounded with
$\norm{D_b}<10$; $\norm\kappa_{\Hh}=1$ and $L\kappa=-\kappa/2$; and
\[
 Cu=u-2m(u)\kappa
\]
is the orthogonal projection reducing $L$ to $L_c$.  These are the actual
form-realization and normalization results of the main text, not consequences
claimed for an arbitrary computable quadratic form.  No identification with the
closure of a separately specified minimal symmetric operator is needed here.
Lemma~\ref{a:lem:radical} supplies independent vectors
\[
 r_j=\kappa^{(2j)}-4^{-j}\kappa\in X_c,
 \qquad q_0(r_j,z)=0\quad(z\in X).
\]
Consequently $r_j\in\ker L_c$.  Only
$\Nn:=\overline{\operatorname{span}\{r_j:j\ge1\}}^{\Hh}
\subseteq\ker L_c$ is used; equality is not assumed.  If
\begin{equation}\label{vfourC:eq:mzero}
 m_0=\inf\sigma(L_c),
\end{equation}
then the correction bound and a nonzero $r_j$ give the safe initial enclosure
$-16\le m_0\le0$.

\begin{theorem}[Conditional computability of the actual centered energy]
\label{vfourC:thm:computable}
Under the stated complete form realization and the effective arithmetic
estimates specified below, a recursive procedure takes any rational
$\varepsilon>0$ and terminates with rational endpoints
\[
 \ell\le m_0\le u,\qquad u-\ell\le\varepsilon.
\]
Any fixed finite independent derivative-radical bank may be imposed exactly,
including the empty bank.  The procedure requires neither a projection onto the
entire kernel nor a convergence rate for such finite banks.
\end{theorem}

The proof is a construction of terminating searches with finite rational
certificates.  It is not a claim that the complete algorithm has been
implemented, that a certificate for the actual sign has been evaluated, or that
a useful running-time bound is known.

\subsection{Effective arithmetic and a rational form core}
\label{vfourC:subsec:arithmetic}

Here a computable scalar or integral means that arbitrarily narrow rational
intervals enclosing it can be produced.  All uses of an upper bound in a finite
certificate mean a rational outward upper enclosure.  The theta series and each
fixed derivative have effective Gaussian remainders on compact intervals.  For
$x\ge0$ one may use
\begin{equation}\label{vfourC:eq:theta}
 W(x)\ge\pi^2 e^{4x-\pi e^{2x}},\qquad
 |\kappa^{(s)}(x)|\le A_s e^{(2s+9/2)x-\pi e^{2x}},
\end{equation}
where a rational $A_s$ can be obtained effectively for every fixed $s$.
Indeed, the first positive theta summand and $\pi>3$ give the lower bound.
For the upper bound use the polynomial recursion in the proof of
Lemma~\ref{a:lem:radical}, bound its finitely many coefficients, and sum a finite
Gaussian head with a geometric bound for its remainder.  Reflection handles
negative $x$.  These bounds give explicit integrable tails for weighted radical
norms, moments, and Fourier transforms on compact frequency intervals.  On a
compact physical interval, a positive lower bound for $W$ is effectively
certifiable by subdivision and a computable derivative bound.  Thus integration
against $1/W$ on that interval is effective as well.

The complete Gamma multiplier has the representation
\begin{equation}\label{vfourC:eq:gamma}
 \begin{split}
 d_\Gamma(t)&=\sum_{h\ge0}\frac{2t^2}{a_h(a_h^2+t^2)},
       \qquad a_h=2h+\tfrac12,\\
 0\le d_\Gamma(t)-\sum_{h<L}\frac{2t^2}{a_h(a_h^2+t^2)}
   &\le2T^2\left(a_L^{-3}+\frac1{4a_L^2}\right)
       \qquad(|t|\le T).
 \end{split}
\end{equation}
In addition,
\begin{equation}\label{vfourC:eq:gamma-bounds}
 0\le d_\Gamma(t)\le17t^2,\qquad
 \operatorname{Lip}(d_\Gamma|_{[-T,T]})\le34T,\qquad
 d_\Gamma(t)\ge\tfrac12\log(1+4t^2).
\end{equation}
The series follows by expanding the Gamma difference kernel into exponentials;
the remainder and upper bounds follow by comparison with the sum of $a_h^{-3}$.
For the lower bound compare the positive decreasing summand, viewed as a
function of $a_h$, with its integral on $[1/2,\infty)$ with spacing two.
The corresponding digamma formulas are consistent with
\cite[Eq.~5.9.16 and Eq.~5.11.2]{DLMF}.  None of these estimates replaces the
complete Gamma form by a truncated form.

For the complete omitted prime form we may use the already rational bound of
Lemma~\ref{a:lem:prime}:
\begin{equation}\label{vfourC:eq:prime-tail}
 |p[u]-p_N[u]|\le\bar\tau_N\norm u_{\Hh}^2,\qquad
 \bar\tau_N=96(N+1)^3(27/512)^{N+1}\longrightarrow0
 \quad(N\ge2).
\end{equation}
This effective cap suffices in place of the sharper exponential cap in
\cite{vfourComputable}.  Every finite $\Lambda(n)$ is computable by finite
factorization.  Both orientations and all prime powers remain present in
\eqref{vfourC:eq:prime-tail}.

\begin{lemma}[Effectively enumerable ambient and centered form cores]
\label{vfourC:lem:core}
There is a recursive list of even, compactly supported, $C^1$ piecewise cubic
splines with rational knots and coefficients in $\mathbb Q+i\mathbb Q$, whose
values and first derivatives vanish at the support endpoints, which is dense in
$X$.  Its exact images under $C$ are dense in $X_c$.
\end{lemma}
\begin{proof}
Lemma~\ref{a:lem:Xcore} already gives the compact smooth form core, so only
rational density and effective enumeration require proof.  On a symmetric
rational mesh choose rational complex Hermite values and first derivatives at
the nonnegative knots, impose the even reflection rules, and impose zero value
and derivative at the outer endpoints.  At zero impose zero first derivative.
Cubic Hermite interpolation then has rational complex coefficients on each
interval and is even and $C^1$.  Enumerating the finite rational data enumerates
all such splines.

For an even compact smooth function, choose a rational support interval slightly
larger than its support.  Hermite interpolants on increasingly fine symmetric
rational meshes converge in ordinary $H^1$.  On each fixed mesh the interpolation
map from its finite endpoint data to $H^1$ is continuous, so rational
approximations to the values and derivatives preserve convergence; the imposed
outer zero data agree with the target.  On the common compact support,
$1/W$ is bounded, while \eqref{vfourC:eq:gamma-bounds} bounds the Gamma form by
$17\norm{v'}_2^2$.  Hence this approximation converges in $X$.  Finally $m$ is
$\Hh$-bounded and $\kappa\in X$, so $C$ is bounded on $X$ and maps the dense list
to a dense list in $X_c$.  This argument needs no computable approximation rate
for an arbitrary element of $X$.
\end{proof}

Each listed spline $v$ has computable $q_0[v]$, $\norm v_{\Hh}^2$, and $m(v)$.
To see the only unbounded-frequency issue explicitly, set $D_v=\norm{v''}_1$;
two integrations by parts give
\[
 |\widehat v(t)|\le\frac{D_v}{\sqrt{2\pi}\,t^2},\qquad
 \int_{|t|>T}d_\Gamma(t)|\widehat v(t)|^2\,dt
       \le\frac{17D_v^2}{\pi T}.
\]
Here $v''$ is the piecewise ordinary second derivative; the $C^1$ endpoint
conditions eliminate distributional boundary terms.  Compact quadrature and
\eqref{vfourC:eq:gamma} compute the retained Gamma part, while
\eqref{vfourC:eq:prime-tail} pays the entire omitted prime sum.  Exact centering
uses the identities
\begin{equation}\label{vfourC:eq:centering}
 q_0[Cv]=q_0[v]+2|m(v)|^2,\qquad
 \norm{Cv}_{\Hh}^2=\norm v_{\Hh}^2-4|m(v)|^2.
\end{equation}
The moment is a computable real or complex number, not a rounded value declared
to vanish exactly.

For rational $b<0$, enumerate the trials and their interval precision until a
strictly negative rational upper enclosure for
$q_0[Cv]-b\norm{Cv}_{\Hh}^2$ is obtained.  This is a sound semidecision for
$m_0<b$.  It is complete by the variational principle, form-core density, and
continuity of the shifted form.  It need not terminate at $m_0=b$.

\subsection{True finite null Grams and the shifted positive reference}
\label{vfourC:subsec:reference}

Fix a finite independent bank $r_1,\ldots,r_j$ and define the synthesis map
$V_jz=\sum_{i=1}^jz_i r_i$.  Its actual weighted Gram and projection are
\begin{equation}\label{vfourC:eq:null-gram}
 H_j=V_j^*V_j
       =[\ip{r_\ell}{r_i}_{\Hh}]_{i,\ell=1}^j,
 \qquad P_j=V_jH_j^{-1}V_j^*.
\end{equation}
These are not identity-Gram substitutions.  Independence makes $H_j$ positive
definite, and effective entry refinement eventually certifies this fact and
computes its inverse, without a prescribed condition number.  Empty-bank
formulas mean $P_0=0$ and omit the corresponding matrices.  The exact map
$(I-P_j)C$ is form-continuous and maps the spline list densely into
$X\cap S$, where
\[
 S=\operatorname{span}\{\kappa,r_1,\ldots,r_j\}^{\perp}\subset\Hh_c.
\]
If $z=V_j^*Cv$, its trial values can be computed from
\[
 q_0[(I-P_j)Cv]=q_0[Cv],\qquad
 \norm{(I-P_j)Cv}_{\Hh}^2=\norm{Cv}_{\Hh}^2-z^*H_j^{-1}z.
\]
Since $b<0$ and the removed bank is exactly null, a strict negative shifted
trial in $X_c$ implies one in $X\cap S$.  Thus the preceding negative
semidecision may also enforce the finite bank exactly.

Fix a rational shift $\sigma>0$.  Because $L_cP_j=0$,
\begin{equation}\label{vfourC:eq:null-decomposition}
 L_c+\sigma I
   =\sigma P_j+(I-P_j)(L_c+\sigma I)(I-P_j)
\end{equation}
as an operator decomposition, or equivalently as a form decomposition.  For
$N\ge2$ such that $\bar\tau_N<\sigma$, put
\[
 \delta'=\sigma-\bar\tau_N>0,\qquad
 c_N(t)=c_\Gamma+2\sum_{2\le n\le N}\frac{\Lambda(n)}{\sqrt n}
                                  \cos(t\log n),
\]
\[
 a_N=(d_\Gamma-c_N)_+,\qquad b_N=(c_N-d_\Gamma)_+,\qquad
 C_N=|c_\Gamma|+2\sum_{2\le n\le N}\frac{\Lambda(n)}{\sqrt n}.
\]
Here $x_+=\max(x,0)$ is computable without deciding the sign of $x$.
Let $A_{\mathrm{ref}}$ represent the ambient positive closed form
\begin{equation}\label{vfourC:eq:reference-form}
 p_{\mathrm{ref}}[u]=\int_{\R}a_N(t)|\widehat u(t)|^2\,dt
                         +\delta'\norm u_{\Hh}^2,
\end{equation}
and let $A_S$ be its form restriction to $S$.  Since
$|a_N-d_\Gamma|\le C_N$, these forms have domains $X$ and $X\cap S$,
respectively, and $A_S\ge\delta'I$.  The complete tail bound gives
\begin{equation}\label{vfourC:eq:head-comparison}
 q_0[u]+\sigma\norm u_{\Hh}^2
 \ge p_{\mathrm{ref}}[u]-\int_{\R}b_N(t)|\widehat u(t)|^2\,dt
 \quad(u\in X\cap S).
\end{equation}
Choose a rational $\bar C>C_N$ and an integer $T$ for which
$\tfrac12\log(1+4T^2)>\bar C$ is certified.  Such a search terminates, and
\eqref{vfourC:eq:gamma-bounds} implies $b_N=0$ for $|t|\ge T$.
No equality test at a zero of $d_\Gamma-c_N$ occurs.

Write $B_S$ for the nonnegative bounded operator on $S$ with form
$\int b_N|\widehat u|^2$, and set
\begin{equation}\label{vfourC:eq:BS}
 K=A_S^{-1/2}B_SA_S^{-1/2}.
\end{equation}
For later computation let $V=(\kappa,r_1,\ldots,r_j)$ denote the finite synthesis
map, $R=A_{\mathrm{ref}}^{-1}$, and $P_S$ the weighted orthogonal projection onto
$S$.  The inverse restricted by these \emph{exact} finite constraints is
\begin{equation}\label{vfourC:eq:restricted-inverse}
 R_S=R-RV(V^*RV)^{-1}V^*R=A_S^{-1}P_S,
\end{equation}
where the last operator is extended by zero on $S^\perp$.
Indeed, applied to an input $f\in\Hh$, the first expression belongs to $S$,
and its reference-form pairing with every $z\in X\cap S$ is
$\ip{f}{z}_{\Hh}$.
Uniqueness of the positive restricted-form solution proves the formula.
The finite inverse Gram $V^*RV$ is positive definite because $R$ is positive
and injective and the columns of $V$ are independent.  Its entries, like those
of $H_j$, retain the actual weighted inner product.

\subsection{A compact positive Birman--Schwinger gate}
\label{vfourC:subsec:matrix}

For $0\le t\le T$ define
\[
 g_t(x)=\frac{W(x)\cos(tx)}{\sqrt{2\pi}},\qquad
 h_t=A_S^{-1/2}P_Sg_t.
\]
For even $u$, $\widehat u(t)=\ip{u}{g_t}_{\Hh}$.  With the convention
$(hh^*)u=\ip{u}{h}_{\Hh}h$, one therefore has
\begin{equation}\label{vfourC:eq:feature}
 K=2\int_0^T b_N(t)h_th_t^*\,dt.
\end{equation}
The normalization $\int c^2W=1/4$ implies
$\int W\le1/4$ and $\int x^2W\le1/2$; for the second inequality use
$2y^2\le\cosh^2y$ at $y=x/2$.  Hence, with
$G_0=(8\pi)^{-1/2}$ and $G_1=(4\pi)^{-1/2}$,
\begin{equation}\label{vfourC:eq:feature-bounds}
 \norm{h_t}_{\Hh}\le\frac{G_0}{\sqrt{\delta'}},\qquad
 \norm{h_t-h_s}_{\Hh}\le\frac{G_1}{\sqrt{\delta'}}|t-s|.
\end{equation}
In particular the integral in \eqref{vfourC:eq:feature} converges in trace norm,
so $K$ is compact and positive.  Compactness, not just boundedness, will be
essential in the completeness argument.

Let
\[
 L_N=2\sum_{2\le n\le N}\Lambda(n)n^{-1/2}\log n,
 \qquad L_b=34T+L_N.
\]
Then $0\le b_N\le C_N$ and $\operatorname{Lip}b_N\le L_b$.  For $r$ equal
midpoint cells, $\Delta=T/r$ and $t_i=(i+1/2)\Delta$, define
\begin{align}
 K_r&=\sum_{i=0}^{r-1}w_i h_{t_i}h_{t_i}^*,
       \qquad w_i=2b_N(t_i)\Delta,\label{vfourC:eq:finite-K}\\
 M_r&=\operatorname{diag}(\sqrt{w_i})
     [\ip{R_Sg_{t_\ell}}{g_{t_i}}_{\Hh}]_{i,\ell=0}^{r-1}
       \operatorname{diag}(\sqrt{w_i}).\label{vfourC:eq:finite-M}
\end{align}
The nonzero spectra of $M_r$ and $K_r$ coincide.  All displayed features and
operators preserve real functions, so $M_r$ is real symmetric.  Zero weights
cause no problem: the nonnegative square-root map is computable at zero and
no division by a weight is used.  The rank-one Lipschitz estimate and midpoint
integration give
\begin{equation}\label{vfourC:eq:quadrature}
 \norm{K-K_r}\le\eta_{\mathrm{quad}}
 :=\frac{T\Delta}{2\delta'}
       (L_bG_0^2+2C_NG_0G_1).
\end{equation}
For verification replace the constants on the right by rational upper bounds.
The estimate follows by integrating $|t-t_i|$ over each cell and using
$\norm{hh^*-kk^*}\le(\norm h+\norm k)\norm{h-k}$.

Let $\widetilde M$ be a symmetric rational approximation to $M_r$, with a
certified rational bound
$\norm{M_r-\widetilde M}\le\eta_{\mathrm{entries}}$ obtained, for example,
from maximum absolute row sums.  The finite gate is
\begin{equation}\label{vfourC:eq:gate}
 I-\widetilde M
       -(\eta_{\mathrm{entries}}+\eta_{\mathrm{quad}})I>0.
\end{equation}
It is decidable for a given rational matrix by exact Sylvester or
$LDL^{\mathsf T}$ tests.  A successful strict test also supplies a rational
$0<\gamma\le1$ for which $I-K\ge\gamma I$: refine a rational lower margin for
the tested matrix and use the spectral correspondence, allowing for the zero
spectrum of $K_r$.  Consequently \eqref{vfourC:eq:head-comparison} yields
\[
 q_0[u]+\sigma\norm u_{\Hh}^2
        \ge\gamma\delta'\norm u_{\Hh}^2\qquad(u\in X\cap S).
\]
Equation~\eqref{vfourC:eq:null-decomposition} restores the removed radical bank
with margin $\min\{\gamma\delta',\sigma\}>0$ on all of $X_c$.

\subsection{Computable inverse upper bounds and primal lower bounds}
\label{vfourC:subsec:inverse}

We now compute the entries in \eqref{vfourC:eq:restricted-inverse} and
\eqref{vfourC:eq:finite-M}; an inverse oracle is not an assumption.  Let
$\mathcal F_2f=\widehat f$ denote the ordinary unitary Fourier transform, let
$J:\Hh\to\Le$ be inclusion, and fix a rational
$\bar W\ge\norm W_\infty$, for instance $\bar W=2$ using
Lemma~\ref{a:lem:explicitW}.  For integers $k\ge1$ put
\[
 m_k(t)=\min(a_N(t),k)\one_{[-k,k]}(t),\quad
 S_k=J^*\mathcal F_2^{-1}m_k\mathcal F_2J,\quad A_k=\delta'I+S_k.
\]
The nonnegative forms increase to $p_{\mathrm{ref}}$, so monotone convergence
of closed forms gives $A_k^{-1}\downarrow R$ strongly.  Since
$0\le S_k\le k\bar W I$, the rational constants and operators
\[
 \alpha_k=\delta'+k\bar W,\qquad
 Y_k=I-A_k/\alpha_k,\qquad q_k=1-\delta'/\alpha_k
\]
satisfy $0\le Y_k\le q_k I$ with $q_k<1$.  Thus, for $\ell\ge0$,
\begin{equation}\label{vfourC:eq:inverse-upper}
 P_{k,\ell}=\alpha_k^{-1}\sum_{s=0}^{\ell}Y_k^s,
 \qquad
 R\le A_k^{-1}\le P_{k,\ell}
                  +q_k^{\ell+1}\delta'^{-1}I.
\end{equation}
The polynomial $P_{k,\ell}$ by itself is \emph{not} a claimed lower bound for
$R$.

For each finite input vector $f$ needed here, the upper quadratic endpoint in
\eqref{vfourC:eq:inverse-upper} is computable.  Its degree-zero term uses
$\norm f_{\Hh}^2$, and for $n\ge1$ the other terms reduce to compact cube
integrals
\begin{equation}\label{vfourC:eq:cube}
 \begin{split}
 \ip{S_k^nf}{f}_{\Hh}
 ={}&\int_{[-k,k]^n}\overline{\widehat f(t_1)}\widehat f(t_n)
       \prod_{h=1}^{n}m_k(t_h)\\
 &\hspace{13mm}\cdot
       \prod_{h=1}^{n-1}\frac{\widehat W(t_h-t_{h+1})}{\sqrt{2\pi}}
       \,dt_1\cdots dt_n.
 \end{split}
\end{equation}
The family of inputs is $g_{t_i},\kappa,r_1,\ldots,r_j$ and their finite
polarization combinations.  Equation~\eqref{vfourC:eq:theta} gives computable
transforms, moments, weighted norms, and compact-frequency moduli of continuity
for these vectors and for $W$.  On the integration cube,
$\min(a_N,k)$ is computably Lipschitz.  The sharp frequency cutoff may jump at
$\pm k$, so no global Lipschitz assertion about $m_k$ is required: those
endpoints are the boundary of the fixed integration cube.  Subdivision with
these moduli gives certified rational enclosures for \eqref{vfourC:eq:cube}.
The formula uses the full Fourier representation; evenness of $W$ makes the
convolution kernel real and even, consistently with the chosen inner-product
convention.

Independent lower endpoints follow from the variational identity
\begin{equation}\label{vfourC:eq:inverse-lower}
 \begin{split}
 2\operatorname{Re}\ip{v}{f}_{\Hh}-p_{\mathrm{ref}}[v]
       &\le\ip{Rf}{f}_{\Hh},\\
 \sup_{v\text{ in the rational spline list}}
       \bigl(2\operatorname{Re}\ip{v}{f}_{\Hh}
                      -p_{\mathrm{ref}}[v]\bigr)
       &=\ip{Rf}{f}_{\Hh}.
 \end{split}
\end{equation}
The full form is used.  Indeed the gap for an arbitrary $v\in X$ is
$p_{\mathrm{ref}}[v-Rf]$, and the form core approximates $Rf\in X$.
All trial expressions are computable.  In particular, for a spline $v$,
\begin{equation}\label{vfourC:eq:reference-tail}
 \int_{|t|>T}a_N(t)|\widehat v(t)|^2\,dt
       \le\frac{D_v^2}{\pi}
                   \left(\frac{17}{T}+\frac{C_N}{3T^3}\right).
\end{equation}

Dovetail $k,\ell$, the spline trials, and all integration precisions.  Retain the
smallest certified rational upper endpoint from
\eqref{vfourC:eq:inverse-upper} and the largest certified rational lower endpoint
from \eqref{vfourC:eq:inverse-lower}.  For every rational tolerance $\zeta>0$,
strong convergence of $A_k^{-1}$, geometric convergence in $\ell$, and core
density force the rational endpoint gap below $\zeta$ at a finite stage.
The stopping condition is the observable rational comparison
$U_f-L_f<\zeta$, not an unobservable limit or a supplied convergence rate.
Polarization computes all required mixed entries.  Refine the positive matrix
$V^*RV$ until a positive rational lower eigenvalue is certified, compute its
inverse with controlled error, and apply \eqref{vfourC:eq:restricted-inverse}.
This proves termination of every finite inverse-entry job and hence of every
requested matrix-entry enclosure.  No useful bound on their cost follows.

\subsection{Completeness for uniform shifted coercivity}
\label{vfourC:subsec:coercivity}

\begin{proposition}[Sound and complete strict-shift semidecision]
\label{vfourC:prop:strict-shift}
For rational $\sigma>0$, fair enumeration of the finite gates
\eqref{vfourC:eq:gate} is a sound semidecision for
\[
 q_0[u]+\sigma\norm u_{\Hh}^2\ge c\norm u_{\Hh}^2
       \quad(u\in X_c)\quad\text{for some }c>0,
\]
equivalently $m_0+\sigma>0$.  It is complete for this \emph{uniform} coercivity
condition and need not halt when $m_0+\sigma=0$.
\end{proposition}
\begin{proof}
Soundness is the conclusion of the finite gate and the exact radical
decomposition.  For completeness suppose the displayed margin $c>0$ exists,
without requiring that it be known.  Choose $N$ large enough that
$\bar\tau_N<\sigma$ and $2\bar\tau_N<c$.  On $X\cap S$ the reference difference
has the exact expression
\[
 p_{\mathrm{ref}}[u]-\ip{B_Su}{u}_{\Hh}
   =q_0[u]+\sigma\norm u_{\Hh}^2
       +(p-p_N)[u]-\bar\tau_N\norm u_{\Hh}^2,
\]
and is therefore at least $(c-2\bar\tau_N)\norm u_{\Hh}^2$.
If the compact positive operator $K$ had $\norm K\ge1$, it would have a
nonzero eigenvector $y$ at an eigenvalue $\lambda\ge1$.  The nonzero vector
$u=A_S^{-1/2}y\in X\cap S$ would satisfy
\[
 p_{\mathrm{ref}}[u]-\ip{B_Su}{u}_{\Hh}
      =\norm y_{\Hh}^2-\ip{Ky}{y}_{\Hh}
      =(1-\lambda)\norm y_{\Hh}^2\le0,
\]
contradicting the uniform lower bound.  Thus $\norm K<1$.  Increasing midpoint
resolution drives \eqref{vfourC:eq:quadrature} to zero; the terminating inverse
procedures drive every finite matrix-entry error to zero.  Consequently some
finite rational gate succeeds.  Fair dovetailing over eligible $N$, midpoint
resolutions, and finite precision requests reaches it.  Each request involves
only finitely many terminating scalar and inverse jobs, so the search does not
wait forever on one unsatisfied proposed gate.
\end{proof}

Pointwise strict positivity without a uniform weighted-$\Hh$ margin is not the
hypothesis of this completeness statement.  In particular the proof does not
silently replace strict shifted coercivity by nonnegativity at an endpoint.

\subsection{Equality-safe rational brackets and finite obligations}
\label{vfourC:subsec:brackets}

\begin{proof}[Proof of Theorem~\ref{vfourC:thm:computable}]
Start with rational endpoints $[\ell,u]=[-16,0]$.  If $u-\ell>\varepsilon$, set
\[
 a=\ell+(u-\ell)/3,\qquad b=\ell+2(u-\ell)/3.
\]
Then $a<b<0$.  Dovetail the positive gate for the rational shift $\sigma=-a$
and the strict-trial search for $m_0<b$.  If the gate succeeds, replace
$\ell$ by $a$; if the trial succeeds first, replace $u$ by $b$.
Both updates preserve the enclosure.  At least one search terminates:
if $m_0\ge b$, then $m_0-a>0$, so the gate terminates; if $m_0\le a$, then
$m_0<b$, so the trial search terminates; and in the overlap both conditions
hold.  In particular equality at either tested threshold cannot block both
branches.  Each update reduces the width by a factor $2/3$.  Thus for
$0<\varepsilon<16$ at most
\[
 \left\lceil\frac{\log(16/\varepsilon)}{\log(3/2)}\right\rceil
\]
outer updates are needed.  This is only a bound on the number of outer updates,
not a bound on the duration of their inner searches.  For $\varepsilon\ge16$
the initial interval already suffices.
\end{proof}

A finite positive certificate must retain all of the following:
\begin{enumerate}[label=(\roman*)]
 \item the prime cutoff, the complete rational tail cap, the positive rational
       shift $\delta'$, and a certified compact support bound for $b_N$;
 \item the theta and Gamma truncation errors, physical and Fourier tail bounds,
       and rational quadrature enclosures used in all entries;
 \item the actual finite null Gram when the bank is used, the positive inverse
       Gram $V^*RV$, its certified inverse error, and the upper/primal-lower
       endpoint evidence for each required inverse entry;
 \item the rational matrix $\widetilde M$, complete scalar and entry error
       debit, midpoint error debit, and an exact positive-definiteness witness
       for \eqref{vfourC:eq:gate}.
\end{enumerate}
A finite negative certificate consists of the rational spline data, the exact
centering identities (and true finite-bank correction, if imposed), and a
strict rational upper enclosure for the complete shifted trial form.
Verification checks finite evidence; it does not invoke an exact-zero test,
an exact sign oracle, or a computed projection onto the entire kernel.

The conclusion is conditional on the specified actual form realization and
effective arithmetic, as in \cite{vfourComputable}.  It establishes
computability as a real number, not decidability of $m_0=0$.  If $m_0<0$, the
strict negative search eventually supplies a witness; if $m_0=0$, arbitrarily
small absolute brackets alone do not certify equality.  Neither the actual
sign, a practical complexity estimate, an effective infinite-radical repair
rate, nor a proof or disproof of the Riemann hypothesis is supplied.
